\begin{abstract}
Reversible Instance Normalization (RevIN) has become a common default in deep
time-series forecasting, although its gains are not universal. We ask a more
practical question: when is instance normalization preferable to global
standardization? We separate per-window normalization, RevIN's learnable affine
parameters, and normalized-space backpropagation in controlled ablations across
forecasting datasets, horizons, backbones, and random seeds. We evaluate both
future dates and unseen series and report seed uncertainty in addition to mean
error. The study shows how instance normalization trades robustness to offset
and scale heterogeneity against the loss of potentially predictive level and
scale information. This view explains why normalized-space training can help
on heterogeneous collections while global standardization remains competitive
on settings with limited shift, and motivates selecting normalization as an
experimental choice rather than treating RevIN as a universal default.

\keywords{Time series forecasting \and Instance normalization \and
Distribution shift}
\end{abstract}
